% =====================================================================================================================
% Chapter 3: the quantum surprise (zero-point energy), human proofs and Lean side by side.
% Sources: papers/prog_zero_point_optimality, p2_certified_zero_point, p2_quantum_transition,
% prog_superfluid_structure_factor, ar_bogoliubov_sk; Lean formal_cert/Programmes/{ZeroPoint,ZeroPointCert,FeynmanBijl}.
% All numbers are copied from the gen_numbers.tex / gen/*.tex / main.tex files of these papers.
% =====================================================================================================================
\lstdefinestyle{chIIIlean}{language=Lean,
  literate={ℝ}{{$\mathbb{R}$}}1 {ℤ}{{$\mathbb{Z}$}}1 {ℕ}{{$\mathbb{N}$}}1 {≤}{{$\le$}}1 {≥}{{$\ge$}}1 {≠}{{$\ne$}}1
  {→}{{$\to$}}1 {↔}{{$\leftrightarrow$}}1 {∀}{{$\forall$}}1 {∃}{{$\exists$}}1 {∧}{{$\wedge$}}1 {∨}{{$\vee$}}1
  {¬}{{$\neg$}}1 {⊤}{{$\top$}}1 {∈}{{$\in$}}1 {∉}{{$\notin$}}1 {∩}{{$\cap$}}1 {⁻¹}{{$^{-1}$}}1 {√}{{$\surd$}}1
  {π}{{$\pi$}}1 {ρ}{{$\rho$}}1 {κ}{{$\kappa$}}1 {α}{{$\alpha$}}1 {δ}{{$\delta$}}1 {ε}{{$\varepsilon$}}1 {σ}{{$\sigma$}}1
  {•}{{$\bullet$}}1 {∑}{{$\sum$}}1 {²}{{$^2$}}1 {₀}{{$_0$}}1 {₁}{{$_1$}}1 {₂}{{$_2$}}1 {‖}{{$\|$}}1 {⟨}{{$\langle$}}1 {⟩}{{$\rangle$}}1
  {∞}{{$\infty$}}1 {×}{{$\times$}}1 {·}{{$\cdot$}}1 {∘}{{$\circ$}}1
  {ψ}{{$\psi$}}1 {φ}{{$\varphi$}}1 {λ}{{$\lambda$}}1 {μ}{{$\mu$}}1 {θ}{{$\theta$}}1 {τ}{{$\tau$}}1
  {—}{{---}}1 {–}{{--}}1 {é}{{\'e}}1
  {⟪}{{$\langle\!\langle$}}1 {⟫}{{$\rangle\!\rangle$}}1 {ₗ}{{$_\ell$}}1 {ℂ}{{$\mathbb{C}$}}1 {←}{{$\leftarrow$}}1}

\chapter{The quantum surprise: zero-point energy and lattices}\chaptermark{The quantum surprise}\label{ch3:chap}\label{ch3:debut}

\noindent Chapter~1 told a success story: for every \emph{completely monotone} repulsive interaction
$V(r)=g(r^2)$, the triangular lattice minimizes the \emph{static} energy $\tfrac12\sum_{R\neq0}V(|R|)$ among all
density-one configurations (the openai/math upstream statement, \tierA{} conditional on upstream theorems that were
produced by an AI and are not peer reviewed; see also the textbook version~1~\cite{v1}). This chapter asks the following question, a very natural one for a physicist: \emph{what
becomes of this result once quantum mechanics is taken into account?} A real crystal is never at rest. Even at
absolute zero, each atom vibrates about its site, and this vibration costs energy, the \emph{zero-point} energy.

We shall see that the answer is a surprise. For a large family of Gaussians, the triangular lattice is not a local
minimum of the zero-point energy: it is a local \emph{maximum} (\tierX). We then give a \emph{certified} version, in
which the logic and the arithmetic are checked by the Lean kernel (\tierA) and the integrals are enclosed by ball
arithmetic (\tierC). Next we ask whether this produces a genuine phase transition (the answer is negative, within the
self-consistent harmonic approximation). We close with a cousin of the same mathematics: quantum fluids, with the
Feynman--Bijl bound and the Bogoliubov structure factor.

\begin{warnbox}[: what this chapter claims, and what it does not]
All results were produced by AI agents (Claude) under the author's direction, and audited by a separate AI agent.
\textbf{No human has audited them.} Nothing in this chapter depends on a theorem of the openai/math upstream.
Montgomery's theorem on Gaussians, used for the static energy, is cited second hand.
The statements about the zero-point energy are Tier~X or Tier~C. Only the logic, the algebra and the rational
arithmetic are Tier~A. The \tierA{} label always attaches to the Lean statement, never to its physical reading.
\end{warnbox}

%=====================================================================================================================
\section{The harmonic crystal from first principles}\label{ch3:sec:cristal}

\subsection{The quantum harmonic oscillator and its minimal energy}\label{ch3:sec:oscillateur}
A particle of mass $m$ in the potential $\tfrac12m\omega^2x^2$ is described by a wave function
$\psi:\R\to\mathbb{C}$ of norm $\int|\psi|^2=1$. Its mean energy is
\[
\langle H\rangle_\psi=\frac{\hbar^2}{2m}\int_\R|\psi'(x)|^2\,dx+\frac{m\omega^2}{2}\int_\R x^2|\psi(x)|^2\,dx .
\]
In classical mechanics the minimal energy is $0$ (the particle at rest at the bottom of the well). In quantum
mechanics this is no longer possible.

\begin{proposition}[Zero-point energy of the oscillator]\label{ch3:prop:osc}
For every normalized $C^1$ function $\psi$ such that $x\psi$ and $\psi'$ are square integrable and
$x|\psi(x)|^2\to0$ as $x\to\pm\infty$, one has $\langle H\rangle_\psi\ge\tfrac12\hbar\omega$, with equality if and only if
$\psi(x)=c\,e^{-m\omega x^2/2\hbar}$ with $|c|$ fixed by the normalization. \tierC
\end{proposition}
\begin{humanproof}
Set $\phi=m\omega\,x\psi+\hbar\psi'$. Expanding $\int|\phi|^2\ge0$:
\[
0\le\int|\phi|^2=m^2\omega^2\!\int x^2|\psi|^2+\hbar^2\!\int|\psi'|^2+2m\omega\hbar\,\mathrm{Re}\!\int x\,\overline{\psi}\,\psi' .
\]
Now $2\,\mathrm{Re}(\overline\psi\psi')=(|\psi|^2)'$, so integrating by parts (the boundary term
$[x|\psi|^2]$ vanishes by hypothesis)
$2\,\mathrm{Re}\int x\overline\psi\psi'=\int x(|\psi|^2)'=-\int|\psi|^2=-1$. What remains is
\[
0\le m^2\omega^2\!\int x^2|\psi|^2+\hbar^2\!\int|\psi'|^2-m\omega\hbar=2m\langle H\rangle_\psi-m\omega\hbar ,
\]
that is, $\langle H\rangle_\psi\ge\hbar\omega/2$. Equality forces $\phi=0$, i.e.\ the linear differential equation
$\psi'=-(m\omega/\hbar)x\psi$, whose solutions are exactly $c\,e^{-m\omega x^2/2\hbar}$.
\end{humanproof}
\noindent The energy $\tfrac12\hbar\omega$ is the \emph{zero-point energy}. Physically, confining the particle
(small $\int x^2|\psi|^2$) forces $\psi$ to vary rapidly (large $\int|\psi'|^2$): this is the uncertainty principle.

\paragraph{Several independent modes \tierC.} If a system has $n$ coordinates $q_1,\dots,q_n$ and a Hamiltonian
$H=\sum_j\bigl(\tfrac{p_j^2}{2m}+\tfrac12m\omega_j^2q_j^2\bigr)$ (\emph{decoupled} modes), the product wave function
$\prod_j\psi_j(q_j)$ shows that the ground-state energy is $\sum_j\tfrac12\hbar\omega_j$. For a crystal, the first
step is therefore to find the coordinates that decouple the vibrations: these are the \emph{phonons}.

\subsection{Warm-up: the linear chain}\label{ch3:sec:chaine}
Take $N$ masses $m$ on a circle, with equilibrium positions $na$, joined to their neighbours by springs of stiffness
$K$, and let $u_n$ be the displacement of mass $n$ (with $u_{n+N}=u_n$). The potential energy is
$\tfrac K2\sum_n(u_{n+1}-u_n)^2$ and Newton's equation reads $m\ddot u_n=K(u_{n+1}+u_{n-1}-2u_n)$.

We look for solutions $u_n=e^{i(kna-\omega t)}$ with $k\in\tfrac{2\pi}{Na}\Z$ (periodicity condition).
Substituting: $-m\omega^2=K(e^{ika}+e^{-ika}-2)=-2K(1-\cos ka)=-4K\sin^2(ka/2)$, hence
\[
\omega(k)=2\sqrt{K/m}\;\bigl|\sin(ka/2)\bigr| .
\]
The discrete Fourier transform decouples the $N$ modes exactly (it diagonalizes the circulant matrix of the potential
energy), so the zero-point energy per particle is
\[
\frac1N\sum_k\frac{\hbar\omega(k)}2\xrightarrow[N\to\infty]{}\frac{\hbar}{2}\cdot\frac1{2\pi}\int_{-\pi}^{\pi}2\sqrt{K/m}\,
|\sin(\theta/2)|\,d\theta=\frac{2}{\pi}\,\hbar\sqrt{K/m}
\]
(a Riemann sum, $\theta=ka$). Three ideas are worth keeping. \textbf{(i)} The frequencies are the square roots of the
eigenvalues of a $k$-dependent matrix. \textbf{(ii)} As $k\to0$, $\omega\approx a\sqrt{K/m}\,|k|$: this is sound, and
the frequency vanishes like $|k|$ (a conical singularity). \textbf{(iii)} The zero-point energy is an \emph{average over
the Brillouin zone} $\theta\in[-\pi,\pi]$. All of this is \tierC.

\subsection{The planar lattice: dynamical matrix and zero-point energy}\label{ch3:sec:matdyn}
A planar Bravais lattice of covolume one is determined, up to rotation, by its \emph{shape} $\tau=x+iy$ ($y>0$), with
basis $a_1=(1/\sqrt y,0)$, $a_2=(x/\sqrt y,\sqrt y)$. The squared length of the vector $ma_1+na_2$ is
\[
Q_{mn}(\tau)=\frac{(m+nx)^2+n^2y^2}{y}.
\]
The triangular lattice corresponds to $\rho=e^{i\pi/3}=\tfrac12+i\tfrac{\sqrt3}2$, the square lattice to $\tau=i$.

\paragraph{The Hessian of a radial potential.} For $V(R)=g(|R|^2)$ with $R\in\R^2$, one has
$\partial_iV=2g'(|R|^2)R_i$ and then
\begin{equation}\label{ch3:eq:hess}
\mathrm{Hess}\,V(R)=2g'(|R|^2)\,I+4g''(|R|^2)\,RR^{\mathsf T}.
\end{equation}
For the Gaussian $g(u)=e^{-\pi au}$: $g'=-\pi a\,g$, $g''=\pi^2a^2g$, hence
$\mathrm{Hess}\,V(R)=(-2\pi a\,I+4\pi^2a^2RR^{\mathsf T})e^{-\pi a|R|^2}$ \tierC.

\paragraph{The equations of motion.} Place a particle of mass $m=1$ at each point $R$ of the lattice $L$, and let
$u_R$ be its small displacement. To second order, the potential energy is
$\tfrac14\sum_{R\ne R'}(u_R-u_{R'})^{\mathsf T}\mathrm{Hess}\,V(R-R')(u_R-u_{R'})$ and Newton gives
$\ddot u_{0}=-\sum_{R\neq0}\mathrm{Hess}\,V(R)\,(u_0-u_R)$. With the plane wave $u_R=e\,e^{i(\mathbf k\cdot R-\omega t)}$
($e\in\mathbb{C}^2$), one obtains $\omega^2e=\sum_{R\ne0}(1-e^{i\mathbf k\cdot R})\mathrm{Hess}\,V(R)\,e$. The terms $R$
and $-R$ have the same Hessian, so one may symmetrize: $1-e^{i\mathbf k\cdot R}$ becomes $1-\cos\mathbf k\cdot R$.

\begin{definition}[Dynamical matrix, stability, zero-point energy]\label{ch3:def:ezp}
For $\mathbf k\in\R^2$, the \emph{dynamical matrix} is the real symmetric $2\times2$ matrix
\[
D(\mathbf k)=\sum_{R\in L\setminus\{0\}}\bigl(1-\cos\mathbf k\cdot R\bigr)\,\mathrm{Hess}\,V(R).
\]
The lattice is \emph{dynamically stable} if $D(\mathbf k)$ is positive semidefinite for every $\mathbf k$. Its two
frequencies are then $\omega_j(\mathbf k)=\sqrt{\lambda_j(D(\mathbf k))}$, and the \emph{harmonic zero-point energy} per
particle (units $\hbar=m=1$) is
\[
E_{\mathrm{zp}}(\tau)=\frac12\Bigl\langle\omega_1+\omega_2\Bigr\rangle_{\mathrm{ZB}}=\frac12\Bigl\langle
\operatorname{tr}D(\mathbf k)^{1/2}\Bigr\rangle_{\mathrm{ZB}},
\]
where $\langle\cdot\rangle_{\mathrm{ZB}}$ is the average over a cell of the dual lattice (the Brillouin zone). \tierC
\end{definition}
\noindent The formula $\operatorname{tr}D^{1/2}$ is the heart of this chapter: the zero-point energy is the trace of the
\emph{square root} of a lattice sum. Unlike the static energy, it is not a sum over pairs.
For a positive semidefinite $2\times2$ matrix, $\sqrt{\lambda_1}+\sqrt{\lambda_2}=\sqrt{\operatorname{tr}D+2\sqrt{\det D}}$
(square both sides), a formula used by the numerical codes.

\paragraph{A sum rule \tierC.} The average of $\cos\mathbf k\cdot R$ over the zone vanishes for $R\neq0$. Hence
$\langle\operatorname{tr}D\rangle_{\mathrm{ZB}}=\sum_{R\ne0}\Delta V(R)$, and by Cauchy--Schwarz and then Jensen
(Exercise~\ref{ch3:ex:einstein}),
\begin{equation}\label{ch3:eq:einstein}
E_{\mathrm{zp}}\le\sqrt{\langle\operatorname{tr}D\rangle_{\mathrm{ZB}}/2}\qquad\text{(Einstein bound).}
\end{equation}
For the Gaussian, Poisson summation gives
$\langle\operatorname{tr}D\rangle=4\pi a-(4\pi^2/a)\,S(1/a)$ with $S(\alpha)=\sum_{P\in L}|P|^2e^{-\pi\alpha|P|^2}$ \tierC.
Luo and Wei prove that, for $\alpha\ge1$, the minimum of $S(\alpha;\cdot)$ over density-one lattices is attained
at the hexagonal lattice~\cite{ch3-luowei} \tierL; hence, for $a\le1$, the triangular lattice \emph{maximizes}
$\langle\operatorname{tr}D\rangle$ \tierC. An upper bound that is maximal at $\rho$, however, says nothing about the
extrema of $E_{\mathrm{zp}}$ itself.

\begin{uniremark}\tierL{} The Poisson summation formula in $\R^d$ states that for a lattice $L\subset\R^d$ of covolume
$|L|$, with dual lattice $L^*=\{P:\,P\cdot R\in\Z\ \forall R\in L\}$, and a Schwartz function $f$,
$\sum_{R\in L}f(R+x)=|L|^{-1}\sum_{P\in L^*}\hat f(P)\,e^{2i\pi P\cdot x}$, with $\hat f(\xi)=\int f(x)e^{-2i\pi\xi\cdot x}dx$
(standard). Applied to $f=\Delta V$ with $V=e^{-\pi a|x|^2}$, it rewrites the direct sum for
$\langle\operatorname{tr}D\rangle$ as a sum over the reciprocal lattice, which converges faster when $a$ is small; this
is the ``Poisson sum in reciprocal space'' used by the numerical implementations below.
\end{uniremark}

\subsection{The role of the triangular point: symmetry and isotropy}\label{ch3:sec:isotropie}
The energy of a lattice depends only on the lattice up to rotation; it is therefore invariant under $\tau\mapsto\tau+1$
and $\tau\mapsto-1/\tau$, hence under the whole group $SL_2(\Z)$. The element $\gamma(\tau)=(\tau-1)/\tau$ fixes $\rho$,
since $\rho^2=\rho-1$, and its derivative at $\rho$ is $1/\rho^2=-\tfrac12-i\tfrac{\sqrt3}2$: this is the rotation by
$-120^\circ$. If $F$ is $C^2$ and invariant, its gradient at $\rho$ is fixed by this rotation, hence zero, and its
Hessian commutes with it, hence is scalar. Thus
\begin{equation}\label{ch3:eq:courbure}
F(\tau)-F(\rho)=c\,\frac{|\tau-\rho|^2}{y_\rho^2}+O(|\tau-\rho|^3),\qquad c=\tfrac14\Delta F(\rho),\quad
\Delta=y^2(\partial_x^2+\partial_y^2).
\end{equation}
\textbf{The sign of $c$ decides everything:} local minimum if $c>0$, local maximum if $c<0$. The two linear-algebra
facts are checked by the kernel \tierAs; the step ``invariance $\Rightarrow$ zero gradient, scalar Hessian''
(chain rule, regularity of $E_{\mathrm{zp}}$) is \tierC.

\begin{leanbox}[: \lean{formal_cert/Programmes/ZeroPoint/Basic.lean}, lines 91--92, 168--171 and 201]
\begin{lstlisting}[style=chIIIlean]
theorem lap_Q (m n x y : ℝ) (hy : y ≠ 0) :
    y ^ 2 * (2 * n ^ 2 / y + 2 * (m + n * x) ^ 2 / y ^ 3) = 2 * Q m n x y := by
theorem rot120_comm_symm (p q r s : ℝ) (hs : s ^ 2 = 3 / 4)
    (h11 : -1 / 2 * p + -s * q = p * (-1 / 2) + q * s)
    (h12 : -1 / 2 * q + -s * r = p * -s + q * (-1 / 2)) :
    q = 0 ∧ r = p := by
theorem rho_sq : rho ^ 2 = rho - 1 := by
\end{lstlisting}
\end{leanbox}
\noindent Reading: \lean{lap_Q} says that $\Delta Q=2Q$ for the length function (one recognizes
$y^2(Q_{xx}+Q_{yy})$, with $Q_{xx}=2n^2/y$ and $Q_{yy}=2(m+nx)^2/y^3$). \lean{rot120_comm_symm} says that a symmetric
matrix $\bigl(\begin{smallmatrix}p&q\\q&r\end{smallmatrix}\bigr)$ commuting with the rotation by $120^\circ$
(with entries $-\tfrac12$ and $\pm s$, $s^2=\tfrac34$) satisfies $q=0$ and $r=p$. The module contains 19 printed
declarations, all with axioms in $\{$\lean{propext}, \lean{Classical.choice}, \lean{Quot.sound}$\}$, and trap
controls (a \lean{sorry} and a planted axiom) caught by \texttt{\#print axioms}~\cite{ch3-p2zp}. \tierAs

\paragraph{The static energy.} For the Gaussian, Montgomery's theorem~\cite{montgomery} asserts that the triangular
lattice minimizes the static energy $E_{\mathrm{st}}=\tfrac12\sum_{R\ne0}e^{-\pi a|R|^2}$ among lattices of
covolume one, for every $a>0$ (\tierL, cited second hand through Deng--Luo~\cite{ch3-dengluo}). For a potential
of amplitude $\varepsilon$ and range $l$, and a mass $m$, the total harmonic energy is
\begin{equation}\label{ch3:eq:total}
E(\tau)=E_{\mathrm{st}}(\tau)+\hbar_{\mathrm{eff}}\,E_{\mathrm{zp}}(\tau),\qquad\hbar_{\mathrm{eff}}=\frac{\hbar}{\sqrt{m\varepsilon l^2}} ,
\end{equation}
in units of $\varepsilon$ and $l$ \tierC. The parameter $\hbar_{\mathrm{eff}}$ measures ``how quantum'' the system
is.

%=====================================================================================================================
\section{The QU-1 conjecture and its numerical refutation}\label{ch3:sec:qu1}

\subsection{The statement, and the rules of the game}
\begin{definition}[Conjecture QU-1]\label{ch3:def:qu1}
For every completely monotone function $g$, the triangular lattice minimizes $E_{\mathrm{zp}}(\tau)$ among planar
Bravais lattices of covolume one.
\end{definition}
\noindent It was the natural ``quantum companion'' of static universal optimality. It was
\emph{preregistered} with a prior probability of $0.6$ of surviving a numerical falsification gate (G1),
an attempted proof for Gaussians (G2) to be launched \emph{only if} G1 survived~\cite{ch3-p2zp}.
The key rule of the Elenchus method applies here: \emph{try to refute before trying to prove}.

\subsection{The counterexample \tierX}\label{ch3:sec:contreex}
For the Gaussian $g(t)=e^{-\pi t}$ ($a=1$), Table~\ref{ch3:tab:head} compares the triangular lattice and the rhombic
lattice $\tau=0.5+0.966025\,i$.
\begin{table}[ht]\centering\small
\caption{Gaussian $a=1$: zero-point energy (spectral zone rule, quadrature error of
about $4.4\times10^{-15}$), stability margin $\min\lambda/|\mathbf k|^2$ over the zone, and exact long-wavelength
limit \tierX~\cite{ch3-p2zp}.}\label{ch3:tab:head}
\begin{tabular}{lrrrr}\toprule
lattice $\tau$ & $E_{\mathrm{zp}}$ & $\Delta E_{\mathrm{zp}}$ & $\min\lambda/|\mathbf k|^2$ & long wavelength\\\midrule
triangular $\rho$ & 1.5033523522 & 0 & 0.03222 & 0.23826\\
$0.5+0.966025\,i$ & 1.4902783183 & $-1.3074\times10^{-2}$ & 0.03493 & 0.10466\\
$i$ (square, negative control) & -- & -- & $-0.1475$ (unstable) & $-0.1475$\\
\bottomrule\end{tabular}
\end{table}

\noindent The rhombic lattice is dynamically stable with a clear margin, and its zero-point energy is
\emph{lower} than that of the triangular lattice. Four implementations of two formulas (direct sum, and Poisson
sum in reciprocal space) agree; a short independent verification program
(\texttt{human\_check.py}, whose core is reproduced in the paper~\cite{ch3-p2zp}) returns $1.490278318$ and $1.503352352$. Of the 625 shapes sampled for this
potential, 65 are stable, and the 64 stable non-triangular shapes all have lower energy: around
$\rho$, the triangular lattice is a \emph{local maximum} within the stable region \tierX.

\paragraph{Twelve potentials.} The conjecture is numerically refuted for 8 of the 12 completely monotone potentials
tested (Table~\ref{ch3:tab:pot}). The \texttt{riesz:4} case is weak: it rests on a single distant competitor, whose
stability margin is 0.08\,\% of the triangular value, and there the triangular point is a local minimum. The
four survivors are \emph{not} proofs: a sample of 567 shapes and two dense paths are not an exhaustive
search.
\begin{table}[ht]\centering\small
\caption{QU-1 verdicts and nature of the triangular point (sign of $c$, interval over 12 directions) \tierX~\cite{ch3-p2zp}.
\texttt{gauss:$a$} is $e^{-\pi ar^2}$, \texttt{yukawa:$\kappa$} is $e^{-\kappa r}/r$, \texttt{riesz:$s$} is $r^{-s}$,
\texttt{dipolar:$d$} is the dipolar exciton.}\label{ch3:tab:pot}
\begin{tabular}{llr@{\,\ldots\,}ll}\toprule
potential & QU-1 verdict & \multicolumn{2}{c}{$c$ at $\rho$} & nature of $\rho$\\\midrule
\texttt{gauss:0.25} & survives & $+0.19$ & $+0.19$ & min\\
\texttt{gauss:0.5} & \textbf{refuted} & $-0.47$ & $-0.45$ & max\\
\texttt{gauss:1} & \textbf{refuted} & $-0.98$ & $-0.95$ & max\\
\texttt{gauss:2} & survives & $+0.48$ & $+0.51$ & min\\
\texttt{gauss:4} & survives & $+0.31$ & $+0.32$ & min\\
\texttt{yukawa:1} & \textbf{refuted} & $-0.40$ & $-0.38$ & max\\
\texttt{yukawa:2} & \textbf{refuted} & $-0.32$ & $-0.31$ & max\\
\texttt{dipolar:0.1} & \textbf{refuted} & $-0.33$ & $-0.32$ & max\\
\texttt{dipolar:0.5} & \textbf{refuted} & $-0.61$ & $-0.58$ & max\\
\texttt{riesz:3} & \textbf{refuted} & $-0.31$ & $-0.29$ & max\\
\texttt{riesz:4} & \textbf{refuted} (weak) & $+0.34$ & $+0.36$ & min\\
\texttt{riesz:6} & survives & $+3.38$ & $+3.45$ & min\\
\bottomrule\end{tabular}
\end{table}

\subsection{The Gaussian family: an interval of local maxima \tierX}\label{ch3:sec:intervalle}
For $V=e^{-\pi ar^2}$, the curvature $c(a)=\tfrac14\Delta E_{\mathrm{zp}}(\rho)$ was computed for 23 values of
$a\in[0.1, 8]$, and its sign changes were then located by bisection (Table~\ref{ch3:tab:ca}).
\begin{table}[ht]\centering\small
\caption{Curvature $c(a)$ of $E_{\mathrm{zp}}$ at the triangular point, Gaussian $e^{-\pi ar^2}$ (excerpt) \tierX~\cite{ch3-p2zp}.}\label{ch3:tab:ca}
\setlength{\tabcolsep}{3pt}\begin{tabular}{rrrrrrrrrrr}\toprule
$a$ & 0.25 & 0.35 & 0.4 & 0.5 & 0.8 & 1 & 1.2 & 1.4 & 1.6 & 2\\\midrule
$c(a)$ & $+0.1882$ & $+0.0796$ & $-0.0799$ & $-0.4605$ & $-1.0707$ & $-0.9620$ & $-0.6598$ & $-0.3035$ & $+0.0285$ & $+0.4945$\\
\bottomrule\end{tabular}
\end{table}

\begin{proposition}[Numerical result, P2]\label{ch3:prop:intervalle}
For the Gaussian $e^{-\pi ar^2}$, among the values $0.15\le a\le4$ explored, the curvature $c(a)$ is negative exactly
on the numerical interval $a\in(0.37668, 1.58148)$; on this interval, the triangular lattice is a
\emph{strict local maximum} of $E_{\mathrm{zp}}$. \tierX
\end{proposition}
\noindent The endpoints are numerical roots obtained by bisection; the scan does not exclude further sign changes
between its points, and at $c=0$ the second-order test is inconclusive. The restricted version QU-1$'$ (strict local
minimum outside the interval, strict local maximum inside) is a conjecture supported at Tier~X, not proved.

\subsection{The mechanism: concavity of the trace of the square root}\label{ch3:sec:mecanisme}
Why does the analogy with the static case fail? The static energy is \emph{linear} in $V$. A completely monotone
potential is a positive mixture of Gaussians (Bernstein's theorem, Chapter~1), and a property that holds for each
Gaussian passes to the mixture. The zero-point energy, by contrast, contains $\operatorname{tr}\sqrt{D}$: it is
not linear in $V$, and the mixture argument does not apply \tierC.

More precisely, at fixed $\mathbf k$, write $\Phi(D)=\operatorname{tr}D^{1/2}$. Differentiating
$\Phi(D(\tau))$ twice gives two terms:
\[
\Delta E_{\mathrm{zp}}=\underbrace{\tfrac12\Bigl\langle y^2\textstyle\sum_i\operatorname{tr}\bigl(\Phi'(D)\,\partial_i^2D\bigr)\Bigr\rangle}_{\mathrm{Lin}\ \text{(linear response)}}
+\underbrace{\tfrac12\Bigl\langle y^2\textstyle\sum_i\Phi''(D)[\partial_iD,\partial_iD]\Bigr\rangle}_{\mathrm{Conc}\ \le\ 0} .
\]
The second term is nonpositive by \emph{concavity} of the square root on positive semidefinite matrices (the matrix
analogue of the concavity of $\sqrt{\cdot}$ on $\R_+$) \tierC. The first has no controlled sign. For $a=1$, the
numerical decomposition gives $\Delta E_{\mathrm{zp}}(\rho)=-3.848=-0.799+(-3.048)$: concavity dominates \tierX. For
$a=0.25$ and $a=2$, on the contrary, a large term $\mathrm{Lin}>0$ wins.

\begin{physbox}[: where is the energy gained?]
For the rhombic competitor at $a=1$, the gain comes mainly from a \emph{softening of the lower branch near the zone
boundary}. The mean squared frequency decreases only slightly: the Einstein bound~\eqref{ch3:eq:einstein} explains only
24\,\% of the gain; the remaining 76\,\% come from the Jensen gap, that is, from a more spread-out spectrum. Since
$\sqrt{\cdot}$ is concave, spreading the spectrum at fixed quadratic mean lowers the mean frequency.
These numbers are \tierX, and their reading is \tierC.
\end{physbox}

\begin{warnbox}[: the competitors are not free of constraints]
The triangular lattice has the lowest static energy; a rhombic competitor is therefore held in its shape only by a
constraint. QU-1 compares a functional over imposed shapes, and it is refuted as such. The physical question
(total energy, each lattice at its equilibrium shape) is that of Section~\ref{ch3:sec:transition}.
\end{warnbox}

%=====================================================================================================================
\section{B2: a certified result}\label{ch3:sec:b2}
A floating-point computation, however robust, proves nothing. The B2 programme sought a \emph{certified} statement at
$a=1$~\cite{ch3-b2}. The preregistered statement asked for a certified two-sided enclosure of the Laplacian
$\Delta E_{\mathrm{zp}}(\rho)<0$, verified in Lean. \textbf{This is not what was obtained.} What was obtained is a
statement that is weaker in one respect (a one-sided, local bound along a single path) and different in another (it uses
neither isotropy nor criticality).

\paragraph{The setting.} Start from the triangular lattice $L_\rho$ (one vector on the $x$-axis) and shear it by the
symmetric matrix of determinant one
\[
P_t=\begin{pmatrix}1&t\\t&1+t^2\end{pmatrix},\qquad L_t=P_tL_\rho,\qquad
E(t)=\frac12\int_{[-\frac12,\frac12]^2}\operatorname{tr}D_t(\kappa)^{1/2}\,d\kappa ,
\]
with $V=e^{-\pi r^2}$, and with the lattice labels and the variable $\kappa$ (reduced coordinate) held fixed as $t$
varies. The shape of the lattice $L_t$ is $\tau(t)=\rho+\sqrt3\,t+O(t^2)$ \tierC.

\begin{theorem}[B2]\label{ch3:thm:b2}
For every $t$ with $0<|t|\le10^{-5}$, one has $E(t)<E(0)$. Consequently the triangular lattice is not a local minimum
of the harmonic zero-point energy for $e^{-\pi r^2}$. \tierC{} overall; logic and arithmetic \tierAs{}; integrals
\tierC{} (ball arithmetic).
\end{theorem}
\noindent The proof consists of five building blocks. We give them in order, each with its tier.

\subsection{Block 1: a trace inequality for \texorpdfstring{$2\times2$}{2x2} matrices}\label{ch3:sec:trace}
\begin{lemma}[Concavity of $\operatorname{tr}\sqrt{\cdot}$ in tangent form]\label{ch3:lem:trace}
Let $X$ be real symmetric positive definite, $S$ real symmetric, and $P=X^{-1}$. Then
\[
2\operatorname{tr}S\le\operatorname{tr}X+\operatorname{tr}(X^{-1}S^2),\qquad\text{with equality if }S=X.
\]
\tierAs{} for $2\times2$ matrices (the only case used); \tierC{} in arbitrary size.
\end{lemma}
\begin{humanproof}
\emph{Scalar case first.} For $x>0$ and $s\in\R$: $x+s^2/x-2s=(s-x)^2/x\ge0$. This is exactly the AM--GM
inequality, and it is also the tangent line to the concave function $\sqrt{\cdot}$.

\emph{Matrix case.} Set $M=S-X$, which is symmetric. Expanding and using $PX=XP=I$:
\[
\operatorname{tr}(PM^2)=\operatorname{tr}(PS^2)-\operatorname{tr}(PSX)-\operatorname{tr}(PXS)+\operatorname{tr}(PX^2)
=\operatorname{tr}(PS^2)-2\operatorname{tr}S+\operatorname{tr}X ,
\]
since $\operatorname{tr}(PSX)=\operatorname{tr}(XPS)=\operatorname{tr}S$ (cyclic invariance of the trace),
$\operatorname{tr}(PXS)=\operatorname{tr}S$ and $PX^2=X$. It therefore suffices to show
$\operatorname{tr}(PM^2)\ge0$. If $m_1,m_2$ are the columns of $M$, then $M^2=MM^{\mathsf T}$ and
$\operatorname{tr}(PMM^{\mathsf T})=\operatorname{tr}(M^{\mathsf T}PM)=m_1^{\mathsf T}Pm_1+m_2^{\mathsf T}Pm_2\ge0$, since
$P$ is positive definite. For equality, $S=X$ gives $M=0$.

\emph{Positivity of the quadratic form, by hand.} If
$P=\bigl(\begin{smallmatrix}p_{11}&p_{12}\\p_{12}&p_{22}\end{smallmatrix}\bigr)$ with $p_{11}>0$ and $\det P\ge0$, then
\[
p_{11}\,(p_{11}a^2+2p_{12}ab+p_{22}b^2)=(p_{11}a+p_{12}b)^2+(p_{11}p_{22}-p_{12}^2)\,b^2\ge0,
\]
and one divides by $p_{11}>0$.
\end{humanproof}

\begin{leanbox}[: \lean{formal_cert/Programmes/ZeroPointCert/Basic.lean}, lines 49--50 and 64--78]
\begin{lstlisting}[style=chIIIlean]
theorem quad_nonneg (p11 p12 p22 a b : ℝ) (hp : 0 < p11) (hdet : 0 ≤ p11 * p22 - p12 ^ 2) :
    0 ≤ p11 * a ^ 2 + 2 * p12 * a * b + p22 * b ^ 2 := by

theorem two_trace_le_of_inverse (x11 x12 x22 p11 p12 p22 s11 s12 s22 : ℝ)
    (hp : 0 < p11) (hdet : 0 ≤ p11 * p22 - p12 ^ 2)
    (h11 : x11 * p11 + x12 * p12 = 1) (h12 : x11 * p12 + x12 * p22 = 0)
    (h21 : x12 * p11 + x22 * p12 = 0) (h22 : x12 * p12 + x22 * p22 = 1) :
    2 * (s11 + s22) ≤ (x11 + x22) +
      (p11 * (s11 ^ 2 + s12 ^ 2) + 2 * p12 * (s12 * (s11 + s22)) + p22 * (s12 ^ 2 + s22 ^ 2)) := by
  have key : (x11 + x22) +
      (p11 * (s11 ^ 2 + s12 ^ 2) + 2 * p12 * (s12 * (s11 + s22)) + p22 * (s12 ^ 2 + s22 ^ 2)) - 2 * (s11 + s22)
      = (p11 * (s11 - x11) ^ 2 + 2 * p12 * (s11 - x11) * (s12 - x12) + p22 * (s12 - x12) ^ 2)
        + (p11 * (s12 - x12) ^ 2 + 2 * p12 * (s12 - x12) * (s22 - x22) + p22 * (s22 - x22) ^ 2) := by
    linear_combination (2 * s11 - x11) * h11 + (2 * s12 - x12) * h12 + (2 * s12 - x12) * h21
      + (2 * s22 - x22) * h22
  have q1 := quad_nonneg p11 p12 p22 (s11 - x11) (s12 - x12) hp hdet
  have q2 := quad_nonneg p11 p12 p22 (s12 - x12) (s22 - x22) hp hdet
  linarith
\end{lstlisting}
\end{leanbox}
\noindent \textbf{Side by side.} The four hypotheses \lean{h11}--\lean{h22} are the four entries of $XP=I$. The
right-hand side is $\operatorname{tr}X+\operatorname{tr}(PS^2)$ written entry by entry for symmetric $S$.
The identity \lean{key} is our ``$\operatorname{tr}(PM^2)=\dots$'', written as a sum of two quadratic forms
evaluated on the two columns $(s_{11}-x_{11},s_{12}-x_{12})$ and $(s_{12}-x_{12},s_{22}-x_{22})$ of $M$; the tactic
\lean{linear_combination} checks that it follows from the four equations of $XP=I$ with the indicated coefficients. The
two calls to \lean{quad_nonneg} are the column-by-column positivity. The Lean statement is restricted to
$2\times2$ matrices, which suffices: there are two phonon branches.

\subsection{Block 2: the linear-response majorant \tierC}\label{ch3:sec:majorant}
At each $\kappa\ne0$, set $S=D_t(\kappa)^{1/2}$ (so $S^2=D_t$) and $X=D_0(\kappa)^{1/2}$, which is positive definite
because the triangular lattice is dynamically stable. Lemma~\ref{ch3:lem:trace} gives
\[
\operatorname{tr}D_t^{1/2}\le\tfrac12\bigl(\operatorname{tr}X+\operatorname{tr}(X^{-1}D_t)\bigr).
\]
Integrating: $E(t)\le U(t):=\tfrac14\int\bigl[\operatorname{tr}X+\operatorname{tr}(X^{-1}D_t)\bigr]d\kappa$, and
$U(0)=E(0)$ by the equality case. \textbf{Why this majorant?} $U$ is \emph{affine} in $D_t$: its dependence on $t$
is a single absolutely convergent lattice sum against the fixed matrix $X^{-1}$, with no square root of the
perturbed matrix. The price is that $U$ bounds $E$ only from above. The neglected term,
$\tfrac12\operatorname{tr}(X^{-1}(S-X)^2)\ge0$, can only help the conclusion ``$E$ is smaller''. Lean sees only
the pointwise inequality; the integration, the integrability and the identification $S=D_t^{1/2}$ are human ``glue''.

\subsection{Block 3: the window bound}\label{ch3:sec:fenetre}
Since $U$ is affine in $D_t$, its symmetric second difference is an integral of the second difference of the Hessians
$H_n(t)=\mathrm{Hess}\,V(P_tR_n)$:
\[
U(t)+U(-t)-2U(0)=\tfrac14\int\operatorname{tr}\Bigl(X^{-1}\textstyle\sum_n(1-\cos2\pi\kappa\cdot n)\bigl[H_n(t)+H_n(-t)-2H_n(0)\bigr]\Bigr)d\kappa .
\]
Entry by entry, Taylor's theorem with the Lagrange remainder gives $H_n(t)+H_n(-t)-2H_n(0)=t^2H_n''(\xi)$ with
$|\xi|<t$. Splitting $H_n''(\xi)=H_n''(0)+(H_n''(\xi)-H_n''(0))$, one obtains for $0<t\le t_0=10^{-5}$:
\begin{equation}\label{ch3:eq:fenetre}
\frac{U(t)+U(-t)-2U(0)}{t^2}\le\underbrace{U''(0)}_{\mathrm{Upp0}}+\underbrace{\Bigl(\int\operatorname{tr}X^{-1}\Bigr)}_{\mathrm{trW}}\cdot
\underbrace{\sum_n\sum_{\text{entries}}\sup_{|\xi|\le t_0}\bigl|H_n''(\xi)-H_n''(0)\bigr|}_{\mathrm{esum}} .
\end{equation}
The reduction is \tierC. The derivatives along the path, $q(t)=|P_tR_0|^2$, $q'$, $q''$, and the product rule for
one Hessian entry, are checked by the kernel (\lean{hasDerivAt_qS}, \lean{hasDerivAt_gaussS'},
\lean{hasDerivAt_entry'}) \tierAs.

\paragraph{The numbers \tierC{} (ball arithmetic).} The three quantities in~\eqref{ch3:eq:fenetre} were rigorously
enclosed in Arb ball arithmetic~\cite{ch3-arb} (python-flint 0.6.0, 96 bits, outward rounding): the
zone is cut into four triangles with vertex $\Gamma$ (Duffy change of variables, which removes the conical
singularity), then into 206\,468 cells, with no failure, in 448~s on one core. The tails of the lattice sum are
bounded explicitly (Table~\ref{ch3:tab:boules}). This code is \emph{not} checked by the kernel: an error in the
bounds would not be detected by the Lean theorems.
\begin{table}[ht]\centering\small
\caption{Enclosures certified by ball arithmetic, $a=1$, shear path, $t_0=10^{-5}$
\tierC~\cite{ch3-b2}.}\label{ch3:tab:boules}
\begin{tabular}{lll}\toprule
quantity & enclosure / bound & meaning\\\midrule
$U''(0)=\mathrm{Upp0}$ & $[-1.947713, -1.900168]$ & second derivative of the majorant\\
$\int\operatorname{tr}X^{-1}\,d\kappa$ & $[1.825520, 1.833534]$ & weight of the perturbation\\
$\mathrm{esum}$ & $\le0.0319642$ & tails included\\
$B_{hi}=\mathrm{Upp0}_{hi}+\mathrm{trW}_{hi}\,\mathrm{esum}_{hi}$ & $-1.841560$ & upper bound in~\eqref{ch3:eq:fenetre}\\
$\lambda_{\min}(B_0)$ & $\ge0.918889$ & stability of $L_\rho$ on each cell\\
$\|B_t-B_0\|$ & $\le0.007528$ & margin $0.911361>0$ for $|t|\le10^{-5}$\\
\bottomrule\end{tabular}
\end{table}

\subsection{Block 4: the rational certificate}\label{ch3:sec:certificat}
The upper ends of the balls are copied into Lean as rationals, and the kernel checks their sign.
\begin{leanbox}[: \lean{formal_cert/Programmes/ZeroPointCert/Basic.lean}, lines 127--144 (comments omitted)]
\begin{lstlisting}[style=chIIIlean]
noncomputable def t0 : ℝ := 1 / 100000
noncomputable def Upp0_hi : ℝ := -475041875131 / 250000000000
noncomputable def trW_hi : ℝ := 366706720019 / 200000000000
noncomputable def esum_hi : ℝ := 31964165579 / 1000000000000
noncomputable def lmin_lo : ℝ := 183777872659 / 200000000000
noncomputable def dB_hi : ℝ := 7528036227 / 1000000000000

theorem t0_pos : 0 < t0 := by norm_num [t0]

theorem cert_B_neg : Upp0_hi + trW_hi * esum_hi < 0 := by
  norm_num [Upp0_hi, trW_hi, esum_hi]

theorem cert_stability : 0 < lmin_lo - dB_hi := by
  norm_num [lmin_lo, dB_hi]
\end{lstlisting}
\end{leanbox}
\noindent \textbf{Checking by hand.} $\mathrm{Upp0}_{hi}=-1.900167500524$, $\mathrm{trW}_{hi}=1.833533600095$,
$\mathrm{esum}_{hi}=0.031964165579$; the product is about $0.0586074$ and the sum about $-1.8415601<0$.
\lean{norm_num} performs the same computation in exact rational arithmetic. The LeanFlow tool \lean{lean.certify}
independently verified both inequalities, rewritten in integers \tierAs. What the kernel does \emph{not} check:
that these rationals really bound the integrals (that is the \tierC{} content of Block~3). That they coincide with the
results file is checked by a script.

\subsection{Block 5: the logic, then the assembly}\label{ch3:sec:logique}
\begin{lemma}[A concave majorant excludes a local minimum]\label{ch3:lem:logique}
Let $E,U:\R\to\R$, $t_0>0$ and $B<0$. Assume $E\le U$ on $[-t_0,t_0]$, $U(0)=E(0)$, and
$U(t)+U(-t)-2U(0)\le t^2B$ for $0<t\le t_0$. Then $0$ is not a local minimum of $E$. If moreover $E$ is even,
then $E(t)<E(0)$ for every $0<t\le t_0$. \tierAs
\end{lemma}
\begin{humanproof}
For $0<t\le t_0$: $E(t)+E(-t)\le U(t)+U(-t)\le2U(0)+t^2B<2U(0)=2E(0)$. Hence at least one of the numbers $E(t)$,
$E(-t)$ is strictly less than $E(0)$. If $0$ were a local minimum, there would exist $\varepsilon>0$ such that
$E\ge E(0)$ on $(-\varepsilon,\varepsilon)$; the choice $t=\min(\varepsilon/2,t_0)$ gives a contradiction. If $E$ is
even, $2E(t)<2E(0)$.
\end{humanproof}
\begin{leanbox}[: \lean{formal_cert/Programmes/ZeroPointCert/Basic.lean}, lines 100--102 and 157--162]
\begin{lstlisting}[style=chIIIlean]
theorem not_isLocalMin_of_majorant {E U : ℝ → ℝ} {t0 B : ℝ} (ht0 : 0 < t0) (hB : B < 0)
    (hEU : ∀ t, |t| ≤ t0 → E t ≤ U t) (h0 : U 0 = E 0)
    (hU : ∀ t, 0 < t → t ≤ t0 → U t + U (-t) - 2 * U 0 ≤ t ^ 2 * B) : ¬ IsLocalMin E 0 := by

theorem triangular_not_localMin_of_certificate (E U : ℝ → ℝ) (Upp0 trW esum : ℝ)
    (hEU : ∀ t, |t| ≤ t0 → E t ≤ U t) (h0 : U 0 = E 0)
    (hwin : ∀ t, 0 < t → t ≤ t0 → U t + U (-t) - 2 * U 0 ≤ t ^ 2 * (Upp0 + trW * esum))
    (hU : Upp0 ≤ Upp0_hi) (hW0 : 0 ≤ trW) (hW : trW ≤ trW_hi) (he0 : 0 ≤ esum) (he : esum ≤ esum_hi)
    (hsym : ∀ t, E (-t) = E t) :
    ¬ IsLocalMin E 0 ∧ ∀ t, 0 < t → t ≤ t0 → E t < E 0 := by
\end{lstlisting}
\end{leanbox}
\noindent \textbf{How to read the assembled theorem.} Each hypothesis is a fact that the analytic work must
supply: \lean{hEU} and \lean{h0} are Block~2, \lean{hwin} is Block~3, \lean{hU}, \lean{hW}, \lean{he} are the
ball enclosures, \lean{hsym} is the mirror symmetry $E(-t)=E(t)$ (the reflection $\mathrm{diag}(1,-1)$ preserves
$L_\rho$ and conjugates $P_t$ into $P_{-t}$, a \tierC{} argument). The Lean theorem is thus the \emph{exact} statement of
what the Tier~C parts must guarantee, not a proof of the physics. The heading of Theorem~\ref{ch3:thm:b2} is \tierC{}
overall, because it rests on Tier~C blocks; only its logic and its arithmetic are Tier~A.

\paragraph{Why this excludes a local minimum of the shape problem \tierC.} Since $\tau(t)=\rho+\sqrt3\,t+O(t^2)$,
small $t\ne0$ give lattices of covolume one, with shape different from $\rho$, arbitrarily close to $\rho$,
and each has strictly lower zero-point energy. A local minimum at $\rho$ would require
$E_{\mathrm{zp}}\ge E_{\mathrm{zp}}(\rho)$ on a neighbourhood, which this path contradicts. Neither criticality nor
isotropy at $\rho$ is used.

\paragraph{Corollary, under additional hypotheses \tierC.} If moreover $E$ is $C^2$, $E_{\mathrm{zp}}$ is $C^2$
near $\rho$, $\rho$ is critical and the Hessian is isotropic (Section~\ref{ch3:sec:isotropie}), then
$E''(0)\le-1.84156$ and $\Delta E_{\mathrm{zp}}(\rho)\le-0.92078<0$. This is compatible with the numerical value
$-3.848$ of P2 \tierX, and much weaker: the method bounds $E''$ by $U''$ from one side only.

\subsection{The controls}\label{ch3:sec:controlesb2}
\begin{itemize}
\item \textbf{Negative control $a=2$} (where P2 finds a local minimum): one obtains $U''(0)\in[+12.69365, +12.75701]$,
the certificate is \emph{not} produced; LeanFlow rejects it (\texttt{false-statement}) and the file \lean{RejectA2.lean}
fails to compile (rc~1). This shows that the method does not always say ``not a minimum''; it does not show
that it is complete. \tierC{} and \tierAs.
\item \textbf{Tampered certificate} (sign of $\mathrm{Upp0}_{hi}$ flipped): \lean{RejectTampered.lean} fails (rc~1) \tierAs.
\item \textbf{Planted \lean{sorry} and axiom}: \lean{Controls.lean} compiles with rc~0, but \texttt{\#print axioms}
displays \lean{sorryAx} and \lean{ProgZeroPointCert.ctl_cheat}. The return code is not the gate; the list of axioms
is \tierAs.
\item \textbf{Non-preregistered by-product, $a=1/2$}: $B_{hi}=-0.918357<0$, same conclusion \tierC.
\end{itemize}

\begin{warnbox}[: exact scope of Tier A here]
The module \lean{ZeroPointCert/Basic.lean} has 18 printed declarations, all with axioms exactly
$\{$\lean{propext}, \lean{Classical.choice}, \lean{Quot.sound}$\}$; it never imports \lean{Mathlib} as a whole.
A search for \lean{skipKernelTC}, \lean{addDecl}, \lean{native_decide}, \lean{implemented_by}, \lean{sorry} and
\lean{axiom} in the files \lean{ZeroPoint/Basic.lean}, \lean{ZeroPointCert/Basic.lean} and
\lean{FeynmanBijl/Basic.lean} finds nothing (checked while this chapter was written). On the other hand, no
\lean{leanchecker} kernel replay (the third check in the definition of Tier~A) is recorded in the folders of these three
modules (P2, B2 and P6). They appear in the replay plan of 2026-10-26 (\texttt{results/kernel\_replay\_2026-10-26/}),
but that replay had not completed for them when this book was assembled: their statements therefore carry the \tierAs{}
label, which rests on compilation by the kernel, the list of axioms and the search above. The eight
\texttt{LEDGER\_UNAUDITED\_TIER\_A} flags of the ledger remain open, because no
person has certified that the Lean statements say what their glosses say.
\end{warnbox}

\paragraph{Work in progress.} Two follow-ups preregistered on 2026-10-26 are \emph{in progress}, and none of their
results is used here: N1 (a classification theorem for Gaussians: is the triangular lattice a local minimum of
$E_{\mathrm{zp}}$ if and only if $a$ lies outside an interval $(a_-,a_+)$?) and N2 (extending the B2 certificate to
physical potentials: Yukawa, dipolar exciton, Riesz).

%=====================================================================================================================
\section{B1: is there a genuine phase transition?}\label{ch3:sec:transition}
The static energy favours the triangular lattice; the zero-point energy (for $a$ in the interval) disfavours it. By
\eqref{ch3:eq:total}, the sum $E_{\mathrm{st}}+\hbar_{\mathrm{eff}}E_{\mathrm{zp}}$ must therefore change minimizer
as $\hbar_{\mathrm{eff}}$ grows. The B1 programme~\cite{ch3-b1} asked: \emph{where, and is it physical?} Its
preregistered stopping rule read: ``if the transition lies beyond melting or disappears beyond the harmonic level,
report it (negative result)''. \textbf{The rule fired: this is a negative result.} B1 contains no
Lean statement.

\subsection{The balance of curvatures}\label{ch3:sec:equilibre}
By~\eqref{ch3:eq:courbure} applied to both energies, the total curvature at $\rho$ is
$c(\hbar_{\mathrm{eff}})=c_{\mathrm{st}}+\hbar_{\mathrm{eff}}\,c_{\mathrm{zp}}$. Montgomery gives $c_{\mathrm{st}}\ge0$,
and $c_{\mathrm{zp}}<0$ on the interval of Proposition~\ref{ch3:prop:intervalle}. The triangular lattice therefore
ceases to be a (harmonic) local minimum at
\[
\hbar^*=-\frac{c_{\mathrm{st}}}{c_{\mathrm{zp}}}\qquad\text{\tierC{} (argument), \tierX{} (numbers).}
\]
At $a=1$: $c_{\mathrm{st}}=0.11913$, $c_{\mathrm{zp}}=-0.96198$, hence $\hbar^*=0.1238$. On $a\in[0.4, 1.5]$,
$\hbar^*$ is minimal ($0.0877$) at $a=0.5$ and diverges at both ends of the interval, where $c_{\mathrm{zp}}\to0$.

\paragraph{A first-order transition \tierC{} and \tierX.} The stabilizer of $\rho$ is the dihedral group $D_3$.
In the disc coordinate $w=(\tau-\rho)/(\tau-\bar\rho)=re^{i\theta}$, an invariant expansion contains $|w|^2$
and $\mathrm{Re}\,w^3$, but not $w^2$: a \emph{cubic term} $r^3\cos3\theta$ is allowed. This is the Landau scenario
of a first-order transition: a state of \emph{finite} distortion catches up with the triangular lattice before the
curvature vanishes. Numerically, at $a=1$, a stable oblique lattice $\tau\approx0.30+0.96\,i$ reaches the triangular
energy already at $\Lambda_1=0.0942<\hbar^*$; the slopes measured along two mirror rays are equal and opposite (ratio
$-1.02$), as the cubic term predicts. The winners lie on the boundary of the harmonic stability region.

\subsection{The Lindemann criterion}\label{ch3:sec:lindemann}
A crystal melts when the root-mean-square amplitude of the vibrations becomes a sizeable fraction of the nearest-neighbour
distance: this is Lindemann's \emph{empirical rule}, not a theorem. In the harmonic theory at $T=0$,
$\langle|u|^2\rangle=\hbar_{\mathrm{eff}}\,U(a)$ with $U=\tfrac12\langle\operatorname{tr}D^{-1/2}\rangle_{\mathrm{ZB}}$
(finite in dimension two, because the $1/|\mathbf k|$ singularity is integrable), and the Lindemann ratio is
$\gamma=\sqrt{\langle|u|^2\rangle}/a_{nn}$ with $a_{nn}^2=2/\sqrt3$. Since $\gamma^2\propto\hbar_{\mathrm{eff}}$, the
estimated melting coupling $\Lambda_L(\gamma_c)=\gamma_c^2a_{nn}^2/U(a)$ is proportional to $\gamma_c^2$: an
error of $\pm10\,\%$ in $\gamma_c$ multiplies it by $0.81$ or $1.21$ \tierC.

Quantum Monte Carlo simulations of two-dimensional bosonic crystals give, verbatim,
``$\gamma=0.230(6)$'' for dipoles, ``$0.235(15)$ for Yukawa bosons'', ``$0.24(1)$ for Coulomb bosons'' and
``a hard-disks $\gamma=0.279$''~\cite{ch3-astra} (\tierC{}, quotation checked verbatim; the citation gate
admits only theorems at Tier~L). Hence the band $\gamma_c\in[0.23, 0.28]$.

\paragraph{Reading \tierX.} At the spinodal point $\hbar^*$, $\gamma\ge0.2798$ for all tested $a$ ($0.313$ at
$a=1$): at the top of the band or beyond. At the first-order point $\Lambda_1$, $\gamma$ ranges from $0.245$ to $0.273$
for $a=0.5$ to $1$: \emph{inside} the band. The Lindemann test alone is therefore \textbf{inconclusive}.

\subsection{The self-consistent harmonic approximation (SCHA)}\label{ch3:sec:scha}
One replaces the bare dynamical matrix by an effective matrix $D_{\mathrm{eff}}$, determined variationally, so that
the potential is averaged over the quantum fluctuations themselves. The relative displacement $w_R=u_R-u_0$ is
then Gaussian, with covariance $C_R=\hbar_{\mathrm{eff}}\langle(1-\cos\mathbf k\cdot R)D_{\mathrm{eff}}^{-1/2}\rangle$.
For a Gaussian potential everything is explicit: with $b=\pi a$, completing the square in the Gaussian integral,
\begin{equation}\label{ch3:eq:lissage}
\bigl\langle e^{-b|R+w|^2}\bigr\rangle_{w\sim\mathcal N(0,C)}=\det(I+2bC)^{-1/2}\exp\!\bigl(-b\,R^{\mathsf T}(I+2bC)^{-1}R\bigr)
\qquad\tierC
\end{equation}
(Exercise~\ref{ch3:ex:gauss}). The fluctuations \emph{smooth} the potential: this is the main anharmonic effect.

\paragraph{Result \tierX.} In the SCHA, at $a=1$ and $\hbar_{\mathrm{eff}}=0.15$, the Laplacian $\Delta E(\rho)$ equals
$+0.117$ (local minimum), whereas the harmonic theory gives $\Delta E(\rho)=4(c_{\mathrm{st}}+\hbar_{\mathrm{eff}}c_{\mathrm{zp}})=-0.101$
(no longer a minimum). The triangular lattice remains a local minimum and the
lowest of the sampled stable lattices up to $\hbar_{\mathrm{eff}}\approx0.27$--$0.33$; there, the SCHA Lindemann
ratio is $0.60$ to $0.87$, far beyond any melting estimate. Two implementations (Python and the
Rust module \texttt{scha.rs} of rusty-SUNDIALS, commit \texttt{9e85e15}) agree to $8.9\times10^{-16}$; a separate AI
auditor reproduced the verdict with its own code and a different quadrature.

\begin{warnbox}[: what the negative result says exactly]
The transition \emph{disappears in the SCHA}. This is not a proof of physical absence: the SCHA is a Gaussian
variational theory that omits the second-order cubic term, and the stiffening it finds is itself of that
order. Path-integral Monte Carlo (PIMC) was \textbf{not} run: \textbf{BLOCKED} (it would require fixed-area box-shape
moves, and finite size masks the signal up to $M\gtrsim12$). Only Bravais lattices were compared. The
interpretation ``ordinary quantum melting comes first'' is an inference from Tier~X numbers, not a statement from the
literature.
\end{warnbox}

\paragraph{Rydberg-dressed atoms are not Gaussian \tierC{} and \tierX.} The interaction of atoms dressed by
a Rydberg state has the form of a ``soft shoulder'' $U_0/(1+(r/R_c)^6)$~\cite{ch3-zeiher,ch3-henkel}. Its planar
Fourier transform takes negative values (relative minimum $-0.0448$ at $kR_c=4.82$, a \tierX{} computation), whereas
that of the Gaussian $e^{-\pi ar^2}$ is $a^{-1}e^{-\pi k^2/a}>0$. Henkel, Nath and Pohl write that, because of the inner
plateau ($r<R_c$), the momentum-space potential has negative contributions around $k_{\mathrm{rot}}\sim2\pi/R_c$, and
they find for this potential a roton-type softening and the formation of a supersolid~\cite{ch3-henkel}.
Prestipino et al.\ write, summarizing the classical literature on cluster crystals, that the Gaussian repulsion is ``a
marginal case''~\cite{ch3-prestipino}. Both quotations were checked verbatim. The more general idea that cluster
crystals require a negative Fourier component is reported only by an unverified secondary reading
and supports no claim here. Rydberg dressing therefore does not realize the Gaussian model of this chapter.

\begin{uniremark}\tierL{} Bochner's theorem (standard): a continuous function $f:\R^d\to\R$ is positive definite, i.e.\
$\sum_{j,l}\bar c_jc_lf(x_j-x_l)\ge0$ for all finite families, if and only if $f$ is the Fourier transform of a finite
positive measure. In particular a radial potential whose planar Fourier transform takes negative values is not positive
definite, whereas the Gaussian $e^{-\pi a|x|^2}$, with transform
$a^{-1}e^{-\pi|k|^2/a}>0$, is; so is every completely monotone $g(|x|^2)$ with finite mixing measure, as a positive
mixture of Gaussians (standard).
\end{uniremark}

%=====================================================================================================================
\section{Quantum fluids: superfluid helium and condensates}\label{ch3:sec:fluides}
A superfluid is not a crystal, but the same mathematics appears in it: an excitation energy, a square root,
and an inequality between a measurable quantity (the static structure factor $S(k)$, obtained by diffraction) and the
spectrum.

\subsection{The Bogoliubov structure factor and its envelope}\label{ch3:sec:bog}
For a dilute Bose gas of mass $m$ and sound velocity $c$, set $h_c=\hbar c$, $k_*=2mc/\hbar$ and
$\varepsilon_0(k)=\hbar^2k^2/2m=h_ck^2/k_*$. The Bogoliubov dispersion relation and the corresponding structure factor
are
\[
\varepsilon_{\mathrm B}(k)=\sqrt{(h_ck)^2+(h_ck^2/k_*)^2},\qquad S_{\mathrm B}(k)=\frac{\varepsilon_0(k)}{\varepsilon_{\mathrm B}(k)} .
\]
In Lean, $S_{\mathrm B}$ is a \emph{definition}; the relation $\varepsilon_0/S_{\mathrm B}=\varepsilon_{\mathrm B}$ is
therefore an algebraic identity, not an independent $N$-body derivation.

\begin{proposition}[Bogoliubov envelope]\label{ch3:prop:enveloppe}
For $h_c,k,k_*>0$: $S_{\mathrm B}(k)^2=k^2/(k_*^2+k^2)$ and $S_{\mathrm B}(k)\le\sqrt{k/2k_*}$, with equality if and
only if $k=k_*$. \tierAs
\end{proposition}
\begin{humanproof}
$S_{\mathrm B}^2=\varepsilon_0^2/\varepsilon_{\mathrm B}^2=\dfrac{h_c^2k^4/k_*^2}{h_c^2k^2+h_c^2k^4/k_*^2}=\dfrac{k^2}{k_*^2+k^2}$
($h_c$ cancels). Both sides being positive,
$S_{\mathrm B}^2\le k/2k_*\iff2k_*k^2\le k(k_*^2+k^2)\iff0\le k(k-k_*)^2$, which is true, with equality if and only
if $k=k_*$. One concludes by taking square roots, which are increasing on $[0,+\infty)$.
\end{humanproof}
\begin{leanbox}[: \lean{formal_cert/Programmes/FeynmanBijl/Basic.lean}, lines 262--268 (comments omitted) and 314--319]
\begin{lstlisting}[style=chIIIlean]
noncomputable def epsFree (hc ks k : ℝ) : ℝ := hc * k ^ 2 / ks
noncomputable def epsBog (hc ks k : ℝ) : ℝ := Real.sqrt ((hc * k) ^ 2 + (hc * k ^ 2 / ks) ^ 2)
noncomputable def sBog (hc ks k : ℝ) : ℝ := epsFree hc ks k / epsBog hc ks k

theorem sBog_le_sqrt {hc ks k : ℝ} (hhc : 0 < hc) (hk : 0 < k) (hks : 0 < ks) :
    sBog hc ks k ≤ Real.sqrt (k / (2 * ks)) := by
  rw [← Real.sqrt_sq (sBog_pos hhc hk hks).le]
  apply Real.sqrt_le_sqrt
  rw [sBog_sq hhc hk hks, div_le_div_iff₀ (by positivity) (by positivity)]
  nlinarith [mul_nonneg hk.le (sq_nonneg (k - ks))]
\end{lstlisting}
\end{leanbox}
\noindent \textbf{Side by side.} The first proof line writes $S_{\mathrm B}=\sqrt{S_{\mathrm B}^2}$, the second
reduces to comparing squares, the third replaces $S_{\mathrm B}^2$ by its closed form and clears denominators, and
the last is our ``$0\le k(k-k_*)^2$'', supplied as a hint to \lean{nlinarith}. This result was first proved
on another Lean stack~\cite{ch3-bsk}, then re-proved on the frozen stack of this book~\cite{ch3-fb}.

\subsection{The Feynman--Bijl bound, in linear algebra}\label{ch3:sec:fb}
The argument of Feynman and Bijl~\cite{ch3-feynman} (these original papers were not read; we follow the formalized
version of~\cite{ch3-fb}) bounds the energy of the lowest density excitation by $\hbar^2k^2/(2mS(k))$.
\paragraph{The finite setting.} Let $E=\mathbb{C}^n$ with the Hermitian inner product $\langle x,y\rangle$ (linear in
$y$), $H$ a Hermitian matrix, $\psi$ an eigenvector, $H\psi=E_0\psi$ with $E_0$ real, and $\rho$ an arbitrary matrix,
with adjoint $\rho^\dagger$. Write $[A,B]=AB-BA$ and define the \emph{excess}
\[
\mathcal A(X)=\langle X\psi,HX\psi\rangle-E_0\langle X\psi,X\psi\rangle=\langle X\psi,(H-E_0)X\psi\rangle .
\]
In physics, $\rho=\hat\rho_{\mathbf k}=\sum_je^{-i\mathbf k\cdot\mathbf r_j}$, $\psi$ is the ground state and
$NS(k)=\|\rho\psi\|^2$.

\begin{lemma}[Double commutator]\label{ch3:lem:dc}
For every matrix $\rho$: $\langle\psi,[\rho^\dagger,[H,\rho]]\psi\rangle=\mathcal A(\rho)+\mathcal A(\rho^\dagger)$. \tierAs
\end{lemma}
\begin{humanproof}
Expand $[\rho^\dagger,[H,\rho]]=\rho^\dagger H\rho-\rho^\dagger\rho H-H\rho\rho^\dagger+\rho H\rho^\dagger$ and
take the expectation of each term in $\psi$.
(a) $\langle\psi,\rho^\dagger H\rho\psi\rangle=\langle\rho\psi,H\rho\psi\rangle$ (definition of the adjoint).
(b) $\langle\psi,\rho^\dagger\rho H\psi\rangle=E_0\|\rho\psi\|^2$ since $H\psi=E_0\psi$.
(c) $\langle\psi,H\rho\rho^\dagger\psi\rangle=\langle H\psi,\rho\rho^\dagger\psi\rangle=E_0\|\rho^\dagger\psi\|^2$
($H$ Hermitian, $E_0$ real).
(d) $\langle\psi,\rho H\rho^\dagger\psi\rangle=\langle\rho^\dagger\psi,H\rho^\dagger\psi\rangle$.
Grouping (a)$-$(b) and (d)$-$(c), one finds $\mathcal A(\rho)+\mathcal A(\rho^\dagger)$.
\end{humanproof}

\begin{lemma}[Rayleigh quotient in a sector]\label{ch3:lem:rayleigh}
Let $(e_i)$ be an orthonormal basis of eigenvectors of $H$, $He_i=\lambda_ie_i$. For every $x\ne0$, there exists $i$ with
$\langle e_i,x\rangle\neq0$ and $\lambda_i\le\langle x,Hx\rangle/\|x\|^2$. \tierAs
\end{lemma}
\begin{humanproof}
With $c_i=\langle e_i,x\rangle$: $\langle x,Hx\rangle=\sum_i|c_i|^2\lambda_i$ and $\|x\|^2=\sum_i|c_i|^2>0$ (Parseval). The
quotient is therefore an average of the $\lambda_i$ weighted by $w_i=|c_i|^2\ge0$. If all the $\lambda_i$ with weight
$w_i>0$ were strictly above the average $\bar\lambda$, one would have $\sum w_i\lambda_i>\bar\lambda\sum w_i=\sum
w_i\lambda_i$, a contradiction.
\end{humanproof}

\begin{theorem}[Feynman--Bijl bound, abstract form]\label{ch3:thm:fb}
Assume $N,m>0$, $S=\|\rho\psi\|^2/N>0$, the f-sum rule
$\langle\psi,[\rho^\dagger,[H,\rho]]\psi\rangle=N\hbar^2k^2/m$ and the symmetry $\mathcal A(\rho^\dagger)=\mathcal
A(\rho)$. Then there exists an eigenvector $e_i$ with $\langle e_i,\rho\psi\rangle\ne0$ and
\[
\lambda_i-E_0\le\frac{\hbar^2k^2}{2mS}. \tierAs
\]
\end{theorem}
\begin{humanproof}
By Lemma~\ref{ch3:lem:dc} and the symmetry, $2\mathcal A(\rho)=N\hbar^2k^2/m$. The trial state
$\phi=\rho\psi/\|\rho\psi\|$ therefore has energy
$\langle\phi,H\phi\rangle=E_0+\mathcal A(\rho)/\|\rho\psi\|^2=E_0+\hbar^2k^2/(2mS)$. Lemma~\ref{ch3:lem:rayleigh}
applied to $x=\rho\psi$ provides $e_i$.
\end{humanproof}
\begin{leanbox}[: \lean{formal_cert/Programmes/FeynmanBijl/Basic.lean}, lines 45--46, 57--58 and 189--194]
\begin{lstlisting}[style=chIIIlean]
noncomputable def excess (H : E →ₗ[ℂ] E) (E0 : ℝ) (ψ : E) (X : E →ₗ[ℂ] E) : ℂ :=
  ⟪X ψ, H (X ψ)⟫_ℂ - (E0 : ℂ) * ⟪X ψ, X ψ⟫_ℂ

theorem double_comm_expect (hH : H.IsSymmetric) (hψ : H ψ = (E0 : ℂ) • ψ) (ρ : E →ₗ[ℂ] E) :
    ⟪ψ, comm (adjoint ρ) (comm H ρ) ψ⟫_ℂ = excess H E0 ψ ρ + excess H E0 ψ (adjoint ρ) := by

theorem feynman_bijl_bound (hH : H.IsSymmetric) (hψ : H ψ = (E0 : ℂ) • ψ) (ρ : E →ₗ[ℂ] E)
    {N m hbar k S : ℝ} (hN : 0 < N) (hm : 0 < m) (hS : 0 < S) (hSdef : S = ‖ρ ψ‖ ^ 2 / N)
    (hf : ⟪ψ, comm (adjoint ρ) (comm H ρ) ψ⟫_ℂ = ((N * hbar ^ 2 * k ^ 2 / m : ℝ) : ℂ))
    (hsym : excess H E0 ψ (adjoint ρ) = excess H E0 ψ ρ) :
    ∃ i, ⟪hH.eigenvectorBasis rfl i, ρ ψ⟫_ℂ ≠ 0 ∧
      hH.eigenvalues rfl i - E0 ≤ hbar ^ 2 * k ^ 2 / (2 * m * S) := by
\end{lstlisting}
\end{leanbox}
\noindent Here \texttt{E} $\to_\ell[\mathbb{C}]$ \texttt{E} denotes the $\mathbb{C}$-linear maps of $E$ (a
finite-dimensional space) and $\langle\!\langle\cdot,\cdot\rangle\!\rangle_{\mathbb{C}}$ the Hermitian inner product.
The four lines \lean{h1}--\lean{h4} of the Lean proof of \lean{double_comm_expect} are exactly (a)--(d). The
Feynman--Bijl module has 21 declarations whose axioms are on the whitelist; it was independently recompiled by a
separate AI auditor, who recovered the same axiom lists and caught the planted controls as well as two mutations of its
own~\cite{ch3-fb}.

\paragraph{Where does the f-sum rule come from? \tierC.} For $N$ bosons with Hamiltonian
$H=\sum_jp_j^2/2m+V(\mathbf r_1,\dots,\mathbf r_N)$, only the kinetic term fails to commute with $\hat\rho_{\mathbf k}$,
and the computation with $[x,p]=i\hbar$ gives the operator identity
$[\hat\rho_{\mathbf k}^\dagger,[H,\hat\rho_{\mathbf k}]]=N\hbar^2k^2/m$.
But the relation $[x,p]=i\hbar$ \emph{has no finite-dimensional realization}: the trace of a commutator vanishes,
whereas that of $i\hbar I$ does not (Exercise~\ref{ch3:ex:trace}). This is why the Lean theorem, set in
finite dimension, takes the f-sum rule as a \emph{hypothesis}. The same holds for the symmetry $\mathbf k\leftrightarrow-\mathbf k$.
A kernel-checked counterexample shows that the symmetry cannot be removed from the halving step
\tierAs, and a kernel-checked witness shows that all the hypotheses are jointly satisfiable and that the bound is
attained \tierAs.

\paragraph{Link with Bogoliubov.} For $S>0$: $\varepsilon_0/S\le\varepsilon_{\mathrm B}\iff S_{\mathrm B}\le S$, and
$\varepsilon_0/S_{\mathrm B}=\varepsilon_{\mathrm B}$ (\lean{feynman_le_epsBog_iff}, \lean{feynman_sBog_eq}) \tierAs.
Hence the \emph{one-sided diffraction test} \tierC: if a measured $S(k)$ exceeds $\sqrt{k/2k_*}$, it is not the
$S_{\mathrm B}$ of a homogeneous Bogoliubov gas with the same sound velocity. A point below the envelope confirms nothing.

\subsection{The test on tabulated data: the roton region}\label{ch3:sec:roton}
\paragraph{Data \tierL.} Table~14.1 of the Donnelly--Barenghi compilation~\cite{ch3-db} (web edition, a
\emph{secondary} compilation, measurements between 1.4 and 4~K, with no printed error bars) gives 99 points $(Q,S(Q))$
after removing $Q=0$ (49 X-ray, 50 neutron). With $c=238.3$~m/s~\cite{ch3-godfrin}, one finds
$k_*=2m_4c/\hbar=3.00379$~\AA$^{-1}$, hence $2k_*=6.008$~\AA$^{-1}$.

\paragraph{Arithmetic \tierX.} With a threshold $\Delta_{ds}=0.082$ fixed \emph{a posteriori} (the gap between the
maxima of the two data sets; it is not in the preregistration), 63 points are refuted, 29 compatible and 7 undecided. In
the roton region $1.7\le Q\le2.4$~\AA$^{-1}$, \textbf{all 21 points are refuted}, with $S$ between $1.62$ and
$2.49$ times the bound. This result does not depend on the threshold: the smallest margin there is $0.328$, about
$4\Delta_{ds}$.

\paragraph{Checking by hand.} At the peak, $Q=2.058$~\AA$^{-1}$ and $S=1.458$, whereas
$\sqrt{Q/2k_*}=\sqrt{2.058/6.008}\approx0.585$. The margin is about $0.87$, well beyond the gap $0.082$
between the data sets.

\paragraph{Reading \tierC.} Liquid helium-4 is not a Bogoliubov gas with its own sound velocity. This is a textbook
fact (helium is not dilute) and \emph{not} a discovery. The reading is conditional on the secondary data,
on temperatures of 1.4 to 4~K being taken as the ground state, and on the a posteriori threshold. For trapped Bose gases,
only values read off figures exist (\tierX): the three absolute points of Stamper-Kurn et al.~\cite{ch3-sk}
lie below a bound averaged over the local density, with low test power, and the points of Steinhauer et
al.~\cite{ch3-st} were multiplied by 2.3 by the authors. No decisive test is possible.

%=====================================================================================================================
\section{Chapter summary}\label{ch3:sec:bilan}
\begin{table}[ht]\centering\small
\caption{The claims of the chapter and their tier.}\label{ch3:tab:bilan}
\begin{tabular}{p{8.6cm}p{5.6cm}}\toprule
claim & tier\\\midrule
$E_{\mathrm{zp}}\ge\tfrac12\hbar\omega$ for the oscillator; phonons of the chain & \tierC\\
identities for $Q$, isotropy ($2\times2$ algebra), $\rho^2=\rho-1$ & \tierAs{} (glue: \tierC)\\
QU-1 refuted for 8 potentials out of 12; local maximum for $a\in(0.37668, 1.58148)$ & \tierX\\
mechanism: concavity of $\operatorname{tr}\sqrt{\cdot}$ & \tierC{} (numbers \tierX)\\
inequality $2\operatorname{tr}S\le\operatorname{tr}X+\operatorname{tr}(X^{-1}S^2)$, logic, rational certificate & \tierAs\\
enclosures of the integrals ($U''(0)\in[-1.948, -1.900]$) & \tierC{} (ball arithmetic)\\
B2: $\rho$ is not a local minimum at $a=1$ & \tierC{} overall\\
B1: no transition in the SCHA; PIMC blocked & \tierX{} (negative)\\
$S_{\mathrm B}\le\sqrt{k/2k_*}$; abstract Feynman--Bijl bound & \tierAs{} (f-sum rule: hypothesis)\\
$^4$He: 21/21 points of the roton region above the envelope & data \tierL, computation \tierX, reading \tierC\\
\bottomrule\end{tabular}
\end{table}
\noindent The methodological lesson is twofold. First, \emph{test before proving}: the QU-1 conjecture, plausible as
it was, could have absorbed months of proof attempts. Second, \emph{separate the layers}: a statement such as
Theorem~\ref{ch3:thm:b2} mixes a kernel-checked core (logic, arithmetic), rigorous but not kernel-checked computation
(balls) and human arguments; the assembled Lean theorem makes this separation explicit, hypothesis by
hypothesis.

%=====================================================================================================================
\section{Exercises}\label{ch3:sec:exercices}

\begin{exercise}[Oscillator]\label{ch3:ex:osc}
(a) Check that $\psi_0(x)=e^{-m\omega x^2/2\hbar}$ satisfies $-\tfrac{\hbar^2}{2m}\psi_0''+\tfrac12m\omega^2x^2\psi_0=\tfrac12\hbar\omega\,\psi_0$.
(b) In the proof of Proposition~\ref{ch3:prop:osc}, justify the integration by parts on $[-A,A]$ and then the
passage to the limit. (c) For the trial state $\psi_\alpha(x)\propto e^{-\alpha x^2}$, compute $\langle H\rangle$ as a
function of $\alpha$ and recover the minimum.
\end{exercise}

\begin{exercise}[Linear chain]\label{ch3:ex:chaine}
(a) Show that the $N\times N$ matrix of the potential energy of the chain is circulant, and that the vectors
$(e^{ikna})_n$, $k\in\tfrac{2\pi}{Na}\Z$, diagonalize it. (b) Establish
$\frac1N\sum_k\tfrac12\hbar\omega(k)\to\tfrac2\pi\hbar\sqrt{K/m}$. (c) Compute the sound velocity $\lim_{k\to0}\omega/|k|$.
\end{exercise}

\begin{exercise}[Radial Hessian]\label{ch3:ex:hess}
(a) Prove formula~\eqref{ch3:eq:hess}. (b) Deduce $\Delta V(R)=4(ug')'(u)$ at $u=|R|^2$, and then, for
$g(u)=e^{-\pi au}$, $\Delta V=(-4\pi a+4\pi^2a^2u)e^{-\pi au}$ (this is the function \lean{m2f} of
\lean{ZeroPoint/Basic.lean}). (c) Why is $D(\mathbf k)$ real symmetric?
\end{exercise}

\begin{exercise}[Einstein bound]\label{ch3:ex:einstein}
(a) Show that $\tfrac12(\omega_1+\omega_2)\le\sqrt{(\omega_1^2+\omega_2^2)/2}$. (b) Using the
Cauchy--Schwarz inequality for the integral over the zone, deduce~\eqref{ch3:eq:einstein}. (c) When does equality hold?
Explain why an upper bound that is maximal at $\rho$ does not, by itself, locate the extrema of $E_{\mathrm{zp}}$.
\end{exercise}

\begin{exercise}[Rotation by $120^\circ$]\label{ch3:ex:rot}
Let $R=\bigl(\begin{smallmatrix}-1/2&-s\\s&-1/2\end{smallmatrix}\bigr)$ with $s^2=3/4$. (a) Show that a vector
$g$ with $Rg=g$ is zero. (b) Show that a symmetric matrix $H$ with $RH=HR$ is scalar (compare with the
statement \lean{rot120_comm_symm}). (c) Show that $\rho^2=\rho-1$ and that $(\tau-1)/\tau$ fixes $\rho$.
\end{exercise}

\begin{exercise}[The trace inequality]\label{ch3:ex:trace2}
(a) Prove Lemma~\ref{ch3:lem:trace} when $X$ and $S$ are diagonal. (b) Give an example of symmetric matrices $X$, $S$,
with $X$ positive definite, for which the inequality is strict, and check it numerically. (c) Show that, for fixed
$X$, $S\mapsto\operatorname{tr}X+\operatorname{tr}(X^{-1}S^2)-2\operatorname{tr}S$ attains its minimum $0$ at
$S=X$. Interpret this as ``the tangent lies above a concave function''.
\end{exercise}

\begin{exercise}[The logic of the majorant]\label{ch3:ex:logique}
(a) Write out the proof of Lemma~\ref{ch3:lem:logique} in full. (b) Construct $E$ and $U$ with $E\le U$, $U(0)=E(0)$
and $U''(0)>0$, but where $0$ is \emph{not} a local minimum of $E$. Deduce that the B2 method can never
certify a local minimum. (c) Why does the hypothesis $E\le U$ not suffice if $U(0)=E(0)$ is replaced by
$U(0)\ge E(0)$?
\end{exercise}

\begin{exercise}[The rational certificate]\label{ch3:ex:cert}
(a) With the rationals of Section~\ref{ch3:sec:certificat}, compute $\mathrm{Upp0}_{hi}+\mathrm{trW}_{hi}\,\mathrm{esum}_{hi}$
to within $10^{-6}$. (b) Up to what value of $\mathrm{esum}_{hi}$ would the certificate remain valid? (Answer:
about $1.0363$.) (c) Prove the monotonicity statement \lean{cert_window}: if $\mathrm{Upp0}\le\mathrm{Upp0}_{hi}$,
$0\le\mathrm{trW}\le\mathrm{trW}_{hi}$ and $0\le\mathrm{esum}\le\mathrm{esum}_{hi}$, then
$\mathrm{Upp0}+\mathrm{trW}\,\mathrm{esum}\le\mathrm{Upp0}_{hi}+\mathrm{trW}_{hi}\,\mathrm{esum}_{hi}$. Where is the
hypothesis $0\le\mathrm{esum}$ used?
\end{exercise}

\begin{exercise}[Gaussian integrals and the SCHA]\label{ch3:ex:gauss}
(a) In dimension one, for $w\sim\mathcal N(0,\sigma^2)$, show that
$\langle e^{-b(R+w)^2}\rangle=(1+2b\sigma^2)^{-1/2}\exp\bigl(-bR^2/(1+2b\sigma^2)\bigr)$ by completing the square.
(b) Deduce~\eqref{ch3:eq:lissage} for a covariance $C=\sigma^2I_2$, and then for an arbitrary symmetric positive
definite $C$ by diagonalizing $C$. (c) Show that the planar Fourier transform of $e^{-\pi ar^2}$ is
$a^{-1}e^{-\pi k^2/a}$ (with the convention $\hat f(k)=\int f(x)e^{-2i\pi k\cdot x}dx$), hence strictly positive.
\end{exercise}

\begin{exercise}[Balance of curvatures and Lindemann]\label{ch3:ex:lind}
(a) With $c_{\mathrm{st}}=0.11913$ and $c_{\mathrm{zp}}=-0.96198$ ($a=1$), recover $\hbar^*\approx0.1238$.
(b) Show that $\Lambda_L(\gamma_c)\propto\gamma_c^2$ and that $\Lambda_1/\Lambda_L=(\gamma(\Lambda_1)/\gamma_c)^2$.
(c) With $\gamma(\Lambda_1)=0.273$ at $a=1$, state for which values of $\gamma_c$ in $[0.23, 0.28]$ the first-order
point precedes the estimated melting. Conclude on the inconclusive nature of the test.
\end{exercise}

\begin{exercise}[Bogoliubov envelope]\label{ch3:ex:bog}
(a) Prove Proposition~\ref{ch3:prop:enveloppe}. (b) Show the phonon regime:
$\frac{k}{k_*+k}<S_{\mathrm B}(k)<\frac{k}{k_*}$. (c) Show that $S_{\mathrm B}(k)\to1$ as $k\to\infty$ and that the free
value $S\equiv1$ violates the envelope for every $0<k<2k_*$: this is the positive control of the test. (d) Redo the
check of the helium-4 peak.
\end{exercise}

\begin{exercise}[Why the f-sum rule is a hypothesis]\label{ch3:ex:trace}
(a) Show that for two square matrices $A,B$ of size $n$, $\operatorname{tr}(AB-BA)=0$. Deduce that there are
no matrices $X,P$ with $XP-PX=i\hbar I_n$. (b) In $\mathbb{C}^2$ with orthonormal basis $(e_0,e_1)$, let
$H=|e_1\rangle\langle e_1|$, $\psi=e_0$ and $\rho=|e_1\rangle\langle e_0|$. Compute $\mathcal A(\rho)$,
$\mathcal A(\rho^\dagger)$ and the double commutator, and check that $\mathcal A(\rho)\ne f/2$: the symmetry is
necessary. (c) With $\rho=|e_1\rangle\langle e_0|+|e_0\rangle\langle e_1|$, $N=\hbar=k=1$ and $m=1/2$, check that the
bound of Theorem~\ref{ch3:thm:fb} is attained.
\end{exercise}

\begin{exercise}[(*) University exercise: Harmonic chain with second neighbours]
Consider the ring of $N$ unit masses of Section~\ref{ch3:sec:chaine} with nearest-neighbour springs of stiffness $K_1>0$
and next-nearest-neighbour springs of stiffness $K_2\ge0$, so that the potential energy is
$\tfrac{K_1}2\sum_n(u_{n+1}-u_n)^2+\tfrac{K_2}2\sum_n(u_{n+2}-u_n)^2$ (lattice spacing $a=1$).
(a) Show that the (here $1\times1$) dynamical matrix is
$D(k)=2K_1(1-\cos k)+2K_2(1-\cos2k)$, i.e.\ $D(k)=\sum_{j\ne0}(1-\cos kj)\,\Phi_j$ with $\Phi_{\pm1}=K_1$,
$\Phi_{\pm2}=K_2$, in analogy with Definition~\ref{ch3:def:ezp}.
(b) Show that $D(k)=4\sin^2(k/2)\,\bigl(K_1+4K_2\cos^2(k/2)\bigr)\ge0$, so the chain is dynamically stable, and that the
sound velocity is $\sqrt{K_1+4K_2}$.
(c) Show that the zero-point energy per particle in the limit $N\to\infty$ is
$E_{\mathrm{zp}}=\frac{\hbar}{2\pi}\int_0^{\pi}2\sin(\theta/2)\sqrt{K_1+4K_2\cos^2(\theta/2)}\,d\theta$, that it reduces
to $\tfrac2\pi\hbar\sqrt{K_1}$ when $K_2=0$, and compute it in closed form for general $K_2$ (substitute
$v=\cos(\theta/2)$).
(d) Verify the Einstein bound~\eqref{ch3:eq:einstein} in this setting, in the form
$E_{\mathrm{zp}}\le\tfrac\hbar2\sqrt{\langle D\rangle}$ with $\langle D\rangle=2K_1+2K_2$, and check it explicitly when $K_2=0$.
\end{exercise}

\begin{exercise}[(*) University exercise: The Gaussian is completely monotone in $r^2$; its Fourier transform]
(a) Show that $g(u)=e^{-\pi au}$ ($a>0$) is completely monotone on $(0,+\infty)$, i.e.\ $(-1)^ng^{(n)}(u)\ge0$ for all
$n\ge0$, and identify its Bernstein measure (a Dirac mass at $t=\pi a$).
(b) Using Fubini and the one-dimensional identity $\int_\R e^{-\pi ax^2-2i\pi\xi x}dx=a^{-1/2}e^{-\pi\xi^2/a}$
(proved by completing the square or by an ODE in $\xi$), show that the Fourier transform on $\R^2$ of
$x\mapsto e^{-\pi a|x|^2}$ is $a^{-1}e^{-\pi|\xi|^2/a}$, consistently with Exercise~\ref{ch3:ex:gauss}(c).
(c) Deduce that for every positive measure $\mu$ on $(0,+\infty)$ with $\int t^{-1}d\mu(t)<\infty$ (which makes $f$
integrable), the radial function
$f(x)=\int e^{-t|x|^2}d\mu(t)$ has a nonnegative Fourier transform on $\R^2$, and give its formula.
\end{exercise}

\label{ch3:fin}
